% ============================================================
% Research-track / derived-conditional block: internal A=B=C gate
% ============================================================
\subsection{Internal Core-Envelope Gate for the GP/NLSE Reduction}
\label{subsec:internal_abc_core_envelope_gate}

\paragraph{Status.}
This subsection is a derived-conditional Research Track result.  It becomes
canon-level only if the resolved SST core is accepted to possess a single
isotropic complex envelope stiffness in the transverse core plane.

Let the resolved core envelope be represented by
\begin{equation}
\Psi(\rho,\theta)=F(\rho)e^{in\theta},
\end{equation}
with winding number \(n\in\mathbb Z\).  The minimal one-modulus isotropic core
energy is
\begin{equation}
E_\perp[\Psi]
=
\kappa\int_{\mathbb R^2}
\left[
|\nabla\Psi|^2
+
\frac{1}{2\xi^2}(1-|\Psi|^2)^2
\right]d^2x .
\label{eq:sst_single_modulus_core_energy}
\end{equation}
In polar coordinates,
\begin{equation}
|\nabla\Psi|^2
=
\left(\frac{dF}{d\rho}\right)^2
+
\frac{n^2F^2}{\rho^2}.
\end{equation}
After setting \(\rho=\xi r\), and dropping the common factor \(2\pi\kappa\),
Eq.~\eqref{eq:sst_single_modulus_core_energy} becomes
\begin{equation}
E_\perp
\propto
\int_0^\infty
\left[
F'^2 r+n^2\frac{F^2}{r}
+\frac12(F^2-1)^2r
\right]dr .
\label{eq:sst_dimensionless_gp_core_energy}
\end{equation}
Comparing this with the general radial form
\begin{equation}
\mathcal L_r
=A F'^2 r+B n^2\frac{F^2}{r}
+\frac{C}{2}(F^2-1)^2r,
\end{equation}
gives
\begin{equation}
\boxed{A=B=C.}
\end{equation}
The Euler--Lagrange equation for the general radial form is
\begin{equation}
F''+\frac{F'}{r}
-\frac{B}{A}n^2\frac{F}{r^2}
+\frac{C}{A}F(1-F^2)=0 .
\end{equation}
Therefore, for the unit vortex sector \(n=1\), the one-modulus condition gives
\begin{equation}
\boxed{
F''+\frac{F'}{r}-\frac{F}{r^2}+F(1-F^2)=0 .
}
\end{equation}
This is the GP/NLSE core equation used in the Track B ring-constant extraction.
The interaction term must consequently be \(\frac12(F^2-1)^2r\).  A coefficient
\(\frac14\) would correspond to \(C/A=1/2\), and would not vary to the unit
GP/NLSE equation above.

Thus the Track B result is canon-derived only under the explicit lemma that the
resolved SST core envelope is governed by the one-modulus isotropic energy
\eqref{eq:sst_single_modulus_core_energy}.  Without this lemma, the GP/NLSE
sector remains a canon-compatible effective-core model rather than a locked
canon theorem.
